\section*{The Appointed Formulary in Full}

Everything the 1962 \latin{Missale Romanum} prints for this Mass, in the order
it prints it: the Latin of the Vatican \latin{editio typica} exactly as the
controlling facsimile has it, each element under its own printed heading,
scriptural reference and marginal number, and beneath it the English of an
identified public-domain witness.

\textbf{Formulary.} \latin{DOMINICA DECIMA QUINTA post Pentecosten},
\latin{II classis}, printed pp.~396--397, marginal nos.\ 1582--1591. It opens in
the lower half of the left column of p.~396, below the Postcommunion of the
Fourteenth Sunday (no.~1581), and ends at the foot of the left column of p.~397,
the right column of that page already carrying the Sixteenth Sunday. It appoints
no Tract, Sequence, blessing or ritual text, and no second or third Collect,
Secret or Postcommunion.

\textbf{The English below.} Scriptural elements are the Douay--Rheims as revised
by Bishop Challoner --- the English of the Clementine Vulgate the missal prints
--- in the Project Gutenberg e-text of that revision, with the psalm loci
resolved in the missal's own Vulgate numbering. The three orations are the
anonymous English of \work{The Roman Missal translated into the English language
for the use of the laity}, first revised edition (Philadelphia: Eugene
Cummiskey, 1861), from the three lines that book supplies for this formulary,
reproduced with its own nineteenth-century spelling and its abbreviated
conclusion ``Thro'.''

\elementguard
\subsection*{1. Introit \cue{Int.}}

\begin{appointedlatin}{\latin{Ant. ad Introitum} --- Ps. 85, 1 et 2-3 --- marginal no.\ 1582}
Inclína, Dómine, aurem tuam ad me, et exáudi me: salvum fac servum tuum, Deus
meus, sperántem in te: miserére mihi, Dómine, quóniam ad te clamávi tota die.
\textnormal{Ps. ibid., 4} Lætífica ánimam servi tui: quia ad te, Dómine, ánimam
meam levávi. ℣.~Glória Patri.
\end{appointedlatin}

\begin{namedtranslation}{English: Douay--Rheims (Challoner), Ps. 85:1--4 printed whole, with the portions the chant sings in bold}
\textbf{v.~1} A prayer for David himself. \textbf{Incline thy ear, O Lord, and
hear me:} for I am needy and poor. \textbf{v.~2} Preserve my soul, for I am
holy: \textbf{save thy servant, O my God, that trusteth in thee.} \textbf{v.~3}
\textbf{Have mercy on me, O Lord, for I have cried to thee all the day.}
\textit{Ps.} \textbf{v.~4} \textbf{Give joy to the soul of thy servant, for to
thee, O Lord, I have lifted up my soul.} ℣.~\latin{Glória Patri.}
\end{namedtranslation}

\englishgap{The antiphon is a three-part cento, not a quotation: v.~1a as far as
\latin{exáudi me}, then v.~2b from \latin{salvum fac}, then v.~3 entire, with the
psalm's own inscription and the halves of vv.~1 and 2 that mention need, poverty
and holiness all unsung. Two of its words are not the Clementine's. \latin{Ad
me} is added at v.~1 and the Douay's ``Incline thy ear, O Lord, and hear me''
has no counterpart for it; \latin{miserére mihi} stands at v.~3 where the Bible
text reads \latin{miserere mei}, and the Douay's ``Have mercy on me'' happens to
answer both. The doxology cue \latin{Glória Patri} belongs to the Ordinary and not
to the proper.}

\elementguard
\subsection*{2. Collect \cue{Coll.}}

\begin{appointedlatin}{\latin{Oratio} --- marginal no.\ 1583}
Ecclésiam tuam, Dómine, miserátio continuáta mundet et múniat: et quia sine te
non potest salva consístere; tuo semper múnere gubernétur. Per Dóminum.
\end{appointedlatin}

\begin{namedtranslation}{English: Cummiskey (Philadelphia, 1861), Collect of this formulary, printed p.~428}
May thy continual mercy purify and defend thy church: and since without thee it
cannot be safe, may it always be directed by the influence of thy grace.
Thro'.
\end{namedtranslation}

\englishgap{\latin{Tuo semper múnere} becomes ``by the influence of thy grace,''
so the noun \latin{munus} --- the gift the Church is governed by, and the word
Augustine uses at the Gradual's psalm to deny that any progress is a merit ---
does not survive into the English, and \latin{semper} is carried by the adverb
of the verb rather than of the gift. The semicolon after \latin{consístere} is
the prayer's hinge; the uncorrected optical layer of the Benziger 1962 shows it
too, and those of the Pustet 1862, the Venice 1570 and the Vatican
\latin{typica} of 1604 show a comma there.}

\elementguard
\subsection*{3. Epistle \cue{Ep.}}

\begin{appointedlatin}{\latin{Léctio Epístolæ beáti Pauli Apóstoli ad Gálatas.} --- Gal. 5, 25-26; 6, 1-10 --- marginal no.\ 1584}
Fratres: Si spíritu vívimus, spíritu et ambulémus. Non efficiámur inánis glóriæ
cúpidi, ínvicem provocántes, ínvicem invidéntes. Fratres, et si præoccupátus
fúerit homo in áliquo delícto, vos, qui spirituáles estis, huiúsmodi instrúite
in spíritu lenitátis, consíderans teípsum, ne et tu tentéris. Alter altérius
ónera portáte, et sic adimplébitis legem Christi. Nam si quis exístimat se
áliquid esse, cum nihil sit, ipse se sedúcit. Opus autem suum probet
unusquísque, et sic in semetípso tantum glóriam habébit, et non in áltero.
Unusquísque enim onus suum portábit. Commúnicet autem is qui catechizátur
verbo, ei qui se catechízat, in ómnibus bonis. Nolíte erráre: Deus non
irridétur. Quæ enim semináverit homo, hæc et metet. Quóniam qui séminat in
carne sua, de carne et metet corruptiónem: qui autem séminat in spíritu, de
spíritu metet vitam ætérnam. Bonum autem faciéntes, non deficiámus: témpore
enim suo metémus, non deficiéntes. Ergo dum tempus habémus, operémur bonum ad
omnes, máxime autem ad domésticos fidei.
\end{appointedlatin}

\begin{namedtranslation}{English: Douay--Rheims (Challoner), Gal. 5:25--26 and 6:1--10, printed whole}
\textbf{5:25} If we live in the Spirit, let us also walk in the Spirit.
\textbf{5:26} Let us not be made desirous of vain glory, provoking one another,
envying one another. \textbf{6:1} Brethren, and if a man be overtaken in any
fault, you, who are spiritual, instruct such a one in the spirit of meekness,
considering thyself, lest thou also be tempted. \textbf{6:2} Bear ye one
another's burdens: and so you shall fulfil the law of Christ. \textbf{6:3} For
if any man think himself to be some thing, whereas he is nothing, he deceiveth
himself. \textbf{6:4} But let every one prove his own work: and so he shall have
glory in himself only and not in another. \textbf{6:5} For every one shall bear
his own burden. \textbf{6:6} And let him that is instructed in the word
communicate to him that instructeth him, in all good things. \textbf{6:7} Be not
deceived: God is not mocked. \textbf{6:8} For what things a man shall sow, those
also shall he reap. For he that soweth in his flesh of the flesh also shall reap
corruption. But he that soweth in the spirit of the spirit shall reap life
everlasting. \textbf{6:9} And in doing good, let us not fail. For in due time we
shall reap, not failing. \textbf{6:10} Therefore, whilst we have time, let us
work good to all men, but especially to those who are of the household of the
faith.
\end{namedtranslation}

\englishgap{The word \latin{Fratres} stands twice and only the first is
liturgical: \latin{Fratres:} at the head is the lesson's added address, prefixed
to Gal. 5:25 and displacing nothing, and it is the missal's own, with no
counterpart in the Douay; \latin{Fratres,} in the middle is Paul's own vocative at
6:1 and the Douay renders it. At 6:6 the Latin transliterates the Greek verb rather than
rendering it, \latin{catechizátur} and \latin{catechízat}, while the Douay reads
``instructed'' and ``instructeth,'' so the Latin and English columns part company
at the one word in this lesson that names a church office. The two burdens of
6:2 and 6:5 are one noun in Latin and in English, \latin{ónera} and
\latin{onus}, where Paul writes two different Greek words. The sowing clause is
numbered differently in the two languages: what the Clementine and the Douay
print at the head of v.~8 stands at the end of Greek v.~7.}

\elementguard
\subsection*{4. Gradual \cue{Grad.}}

\begin{appointedlatin}{\latin{Graduale} --- Ps. 91, 2-3 --- marginal no.\ 1585}
Bonum est confitéri Dómino: et psállere nómini tuo, Altíssime. ℣.~Ad
annuntiándum mane misericórdiam tuam et veritátem tuam per noctem.
\end{appointedlatin}

\begin{namedtranslation}{English: Douay--Rheims (Challoner), Ps. 91:2--3}
\textbf{v.~2} It is good to give praise to the Lord: and to sing to thy name, O
most High. \textbf{v.~3} To shew forth thy mercy in the morning, and thy truth
in the night:
\end{namedtranslation}

\englishgap{The Latin wording is the Clementine's exactly at both verses, with
nothing added, dropped or transposed; what differs is pointing. The Clementine
prints a comma after \latin{misericórdiam tuam} and the typical edition does
not, and the uncorrected optical layers of the Benziger 1962, the Pustet 1862,
the Venice 1570 and the 1604 \latin{typica} all show it. \latin{Confitéri}
carries confession of sin and confession of praise together, and both Augustine
and Cassiodorus build their expositions on the double sense; the Douay's ``to
give praise'' renders the second alone.}

\elementguard
\subsection*{5. Alleluia \cue{All.}}

\begin{appointedlatin}{\latin{Allelúia} --- Ps. 94, 3 --- marginal no.\ 1586}
Allelúia, allelúia. ℣.~\textnormal{Ps. 94, 3} Quóniam Deus magnus Dóminus, et
Rex magnus super omnem terram. Allelúia.
\end{appointedlatin}

\begin{namedtranslation}{English: Douay--Rheims (Challoner), Ps. 94:3 --- the Bible verse the chant departs from}
\textbf{v.~3} For the Lord is a great God, and a great King above all gods.
\end{namedtranslation}

\englishgap{\textbf{The Douay does not answer this chant.} The missal sings
\latin{super omnem terram} where the Clementine reads \latin{super omnes deos},
and the Douay renders the Clementine; the predicate is not an orthographic or
pointing variant but a different one, and it stands already in the uncorrected
optical layers of the Pustet 1862, the Venice 1570 and the 1604 \latin{typica}.
The 1861 Philadelphia hand missal prints its own English of the versicle beside
the Latin at printed p.~428, ``℣.~For the Lord is a great God, and a great King
over all the earth. Alleluia.'' --- an identified nineteenth-century rendering
of what the missal sings, and not the Douay.}

\elementguard
\subsection*{6. Gospel \cue{Gosp.}}

\begin{appointedlatin}{✠ \latin{Sequéntia sancti Evangélii secúndum Lucam.} --- Luc. 7, 11-16 --- marginal no.\ 1587}
In illo témpore: Ibat Iesus in civitátem, quæ vocátur Naim: et ibant cum eo
discípuli eius, et turba copiósa. Cum autem appropinquáret portæ civitátis, ecce
defúnctus efferebátur filius únicus matris suæ: et hæc vídua erat: et turba
civitátis multa cum illa. Quam cum vidísset Dóminus, misericórdia motus super
eam, dixit illi: Noli flere. Et accéssit, et tétigit lóculum. (Hi autem, qui
portábant, stetérunt). Et ait: Aduléscens, tibi dico, surge. Et resédit qui erat
mórtuus, et cœpit loqui. Et dedit illum matri suæ. Accépit autem omnes timor: et
magnificábant Deum, dicéntes: Quia prophéta magnus surréxit in nobis: et quia
Deus visitávit plebem suam.
\end{appointedlatin}

\begin{namedtranslation}{English: Douay--Rheims (Challoner), Lk. 7:11--16}
\textbf{v.~11} And it came to pass afterwards that he went into a city that is
called Naim: and there went with him his disciples and a great multitude.
\textbf{v.~12} And when he came nigh to the gate of the city, behold a dead man
was carried out, the only son of his mother: and she was a widow. And a great
multitude of the city was with her. \textbf{v.~13} Whom when the Lord had seen,
being moved with mercy towards her, he said to her: Weep not. \textbf{v.~14} And
he came near and touched the bier. And they that carried it stood still. And he
said: Young man, I say to thee, arise. \textbf{v.~15} And he that was dead sat up
and began to speak. And he gave him to his mother. \textbf{v.~16} And there came
a fear upon them all: and they glorified God saying: A great prophet is risen up
among us: and, God hath visited his people.
\end{namedtranslation}

\englishgap{The missal's incipit \latin{In illo témpore: Ibat Iesus} replaces
the Bible text's \latin{Et factum est: deinceps ibat}, which the Douay renders
``And it came to pass afterwards that he went''; the missal's opening is its own
and has no Douay counterpart. \latin{Aduléscens} is printed with \emph{u} where
the Clementine reads \latin{Adolescens}, and the uncorrected optical layers of
the Venice 1570 and the 1604 \latin{typica} show the same, and the Douay's
``Young man'' answers the word whichever way it is spelled. The parenthesis at
v.~14 closes before its stop in this book, \latin{stetérunt).}, against the
Clementine's \latin{steterunt.)}.}

\elementguard
\subsection*{7. Offertory \cue{Off.}}

\begin{appointedlatin}{\latin{Ant. ad Offertorium} --- Ps. 39, 2, 3 et 4 --- marginal no.\ 1588}
Exspéctans exspectávi Dóminum, et respéxit me: et exaudívit deprecatiónem meam:
et immísit in os meum cánticum novum, hymnum Deo nostro.
\end{appointedlatin}

\begin{namedtranslation}{English: Douay--Rheims (Challoner), Ps. 39:2--4 printed whole, with the portions the chant sings in bold}
\textbf{v.~2} \textbf{With expectation I have waited for the Lord, and he was
attentive to me.} \textbf{v.~3} And he heard my prayers, and brought me out of
the pit of misery and the mire of dregs. And he set my feet upon a rock, and
directed my steps. \textbf{v.~4} \textbf{And he put a new canticle into my
mouth, a song to our God.} Many shall see, and shall fear: and they shall hope
in the Lord.
\end{namedtranslation}

\englishgap{\textbf{Three of this antiphon's phrases have no Douay equivalent,
because the Douay renders the Clementine and the antiphon does not.} Where the
Bible text has \latin{et intendit mihi}, the chant has \latin{et respéxit me};
where it has \latin{Et exaudivit preces meas}, the chant has \latin{et exaudívit
deprecatiónem meam}; where it has \latin{carmen Deo nostro}, the chant has
\latin{hymnum Deo nostro}. All three stand in the uncorrected optical layers of
the Pustet 1862, the Venice 1570 and the 1604 \latin{typica}. The 1861
Philadelphia hand missal is half a help here, printing at p.~429 ``with
expectation I have waited for the Lord, and he was attentive to me: and he heard
my prayer, and he put a new canticle into my mouth: a hymn to our God'' ---
following the Douay at the first two and the missal at the third. The whole of
v.~3 from \latin{et eduxit me de lacu miseriæ} to \latin{direxit gressus meos}
is cut, and so is the second half of v.~4, which the Douay prints and which the
chant does not reach.}

\elementguard
\subsection*{8. Secret \cue{Sec.}}

\begin{appointedlatin}{\latin{Secreta} --- marginal no.\ 1589}
Tua nos, Dómine, sacraménta custódiant: et contra diabólicos semper tueántur
incúrsus. Per Dóminum nostrum Iesum Christum, Fílium tuum: Qui tecum vivit et
regnat in unitáte.
\end{appointedlatin}

\begin{namedtranslation}{English: Cummiskey (Philadelphia, 1861), Secret of this formulary, printed p.~429}
May thy mysteries, O Lord, preserve us, and always defend us against the attacks
of the devil.
\end{namedtranslation}

\englishgap{The 1861 book prints no Latin incipit cue at this prayer, and no
conclusion: the long form the typical edition sets out to \latin{in unitáte.} is
untranslated there, where the Collect's short \latin{Per Dóminum.} and the
Postcommunion's long form are rendered by the book's abbreviated ``Thro'.'' The
uncorrected optical layers of the Benziger 1962, the Pustet 1862 and the 1604
\latin{typica} all show the short conclusion at this Secret; the long form here
is the Vatican typical edition's own choice at this place and is not a different
prayer.}

\elementguard
\subsection*{9. Communion \cue{Comm.}}

\begin{appointedlatin}{\latin{Ant. ad Communionem} --- Io. 6, 52 --- marginal no.\ 1590}
Panis, quem ego dédero, caro mea est pro sǽculi vita.
\end{appointedlatin}

\begin{namedtranslation}{English: Douay--Rheims (Challoner), Jn. 6:52 printed whole, with the portion the chant sings in bold}
\textbf{v.~52} If any man eat of this bread, he shall live for ever: \textbf{and
the bread that I will give is my flesh, for the life of the world.}
\end{namedtranslation}

\englishgap{\textbf{The Douay answers neither of the antiphon's two departures.}
\latin{Dédero} is a future perfect where the Clementine has the simple future
\latin{dabo}, and the Douay reads ``I will give''; \latin{pro sǽculi vita}
stands where the Clementine has \latin{pro mundi vita}, and the Douay reads
``the life of the world.'' No registered public-domain English renders either
reading distinctly. Both readings stand in the uncorrected optical layers of the
Pustet 1862, the Venice 1570 and the 1604 \latin{typica}. The first half of the
verse is not sung, and \latin{Panis} is capitalised where the Bible text has it
mid-sentence.}

\elementguard
\subsection*{10. Postcommunion \cue{Postcomm.}}

\begin{appointedlatin}{\latin{Postcommunio} --- marginal no.\ 1591}
Mentes nostras et córpora possídeat, quǽsumus, Dómine, doni cæléstis operátio:
ut non noster sensus in nobis, sed iúgiter eius prævéniat efféctus. Per Dóminum
nostrum Iesum Christum, Fílium tuum: Qui tecum vivit et regnat in unitáte.
\end{appointedlatin}

\begin{namedtranslation}{English: Cummiskey (Philadelphia, 1861), Postcommunion of this formulary, printed p.~429}
May the efficacy of these divine mysteries, O Lord, possess both our souls and
bodies: that their effects, not our own sensuality, may always take the lead in
us.  Thro'.
\end{namedtranslation}

\englishgap{\latin{Doni cæléstis operátio} is a singular subject held back to
the end of its clause behind two objects and a vocative, and the English turns
it plural, ``the efficacy of these divine mysteries,'' which also carries
\latin{efféctus} into ``their effects.'' \latin{Mentes} becomes ``souls'' rather
than minds, and \latin{iúgiter} becomes ``always,'' the same adverb the same
book uses for \latin{semper} in the Collect and the Secret. The typical edition
prints no comma after \latin{nostras}, so \latin{Mentes nostras et córpora} is a
single object pair; the uncorrected optical layer of the Benziger 1962 shows a
comma there, and those of the Pustet 1862 and the 1604 \latin{typica} show
none.}

\textbf{Rubrics carried by the formulary.} \latin{II classis} under the heading;
\latin{Credo.} on its own line after the Gospel; \latin{Præfatio de Ss\~{m}a
Trinitate.} immediately after the Secret. The \latin{Gloria} is said and the
Introit carries no \latin{Allelúia}, both by the general rubrics of 1960 --- RG
431~a and RG 429 --- and not by anything printed on these two pages.
